\section{Introduction}
A spatial cue that reliably indicates where a relevant event will occur speeds and sharpens its detection, while a cue that misdirects attention imposes a measurable cost: detection at the uncued location is slower and less sensitive than it would be with no cue at all \citep{posner1980attention,posner1994visual}. This valid-benefit/invalid-cost signature is the canonical operational definition of covert spatial attention and a load-bearing tool across its study \citep{carrasco2011visual}. The cost arises because attention, once committed to the cued location, must be reallocated when the relevant event appears elsewhere, and that reallocation is neither instantaneous nor guaranteed.

How large the invalid-cue cost should be is, on mechanistic grounds, a function of stimulus load. Biased-competition and normalization accounts hold that stimuli within a receptive field or suppressive pool compete, and that attention resolves the competition by biasing the gain of the attended item against its neighbours \citep{desimone1995neural,reynolds1999competitive,reynolds2000attention,reynolds2009normalization}. Under these accounts, adding competitors deepens the suppressive field around the cued item, so that reallocating the top-down bias \emph{away} from a strongly favoured cued location toward an uncued event becomes progressively harder as set size grows. A priority-map view, in which selection reflects a commitment by collicular and cortical priority circuitry to the cued location, predicts the same direction: a more strongly committed map is more costly to override \citep{krauzlis2018selective,krauzlis2013superior,zenon2012attention}. Whether this load-scaling of the invalid-cue cost emerges in a system that learns the task end-to-end, rather than being built in, is open.

We address this in a trained recurrent attention agent that performs a value-directed, spatially cued orientation-change-detection task from raw pixels. A coloured cue marks one location and, by its colour, a reward value \citep{maunsell2004neuronal,luo2019attention,herman2020attention}; a change occurs on half of trials, and the agent declares wait or change each frame. The agent routes detection through a memory pathway whose read-out we can interrogate directly, asking on which location it places its attention as the change unfolds. We compare two stimulus set sizes (four versus nine locations) within the same model, and contrast a cue-gated top-down read-out against a bottom-up variant of the same architecture. The result, previewed here, is that the invalid-cue cost scales catastrophically with set size in the top-down read-out: a recoverable one-frame delay at the low set size becomes a categorical failure to attend the uncued change at the high set size, while the bottom-up read-out is immune. This identifies the load-scaling of the invalid-cue cost as a specific property of cue-gated top-down selection \citep{sprague2018rescuing}, and yields a concrete behavioural prediction.
